% ======================================================================================
% Nom     : reports/latex/sections/05_methodologie.tex
% Rôle    : Section décrivant la méthodologie et le protocole expérimental.
% Auteur  : Maxime BRONNY
% Version : V1
% Cadre   : Réalisé dans le cadre du cours Techniques d'apprentissage artificiel
%           Master 1 Informatique Big Data
% Usage   : Fichier inclus dans main.tex lors de la compilation du rapport.
% ======================================================================================

\section{Méthodologie}
\label{sec:methodologie}

\subsection{Choix des algorithmes}

Nous comparons trois méthodes de classification appartenant à trois familles
distinctes, ce qui donne à la comparaison un intérêt qui dépasse le simple
classement de performances :

\begin{itemize}
  \item la \textbf{régression logistique}, modèle \emph{linéaire}
    probabiliste : c'est la méthode de référence historique du credit scoring
    \cite{hosmer2000logistic}, interprétable (chaque variable reçoit un poids)
    et rapide ;
  \item l'\textbf{arbre de décision CART} \cite{breiman1984cart}, modèle
    \emph{non linéaire} à base de règles : il capture les interactions entre
    variables et produit des décisions lisibles (suites de tests
    « si\ldots{} alors ») ;
  \item les \textbf{$k$ plus proches voisins} \cite{cover1967knn}, méthode
    \emph{non paramétrique} à base d'instances : aucune phase d'apprentissage,
    la prédiction repose entièrement sur la similarité aux exemples connus.
\end{itemize}

Nous avons volontairement limité l'étude à trois méthodes maîtrisées de bout en
bout (toutes trois implémentées from scratch) plutôt que d'en survoler
davantage : l'objectif pédagogique est de comprendre finement chaque algorithme,
ses hyperparamètres et ses modes d'échec.

\subsection{Protocole expérimental}

Le protocole, entièrement automatisé par \texttt{main.py}, est identique pour
les trois modèles :

\begin{enumerate}
  \item \textbf{entraînement} sur les 22\,802 exemples du train ;
  \item \textbf{sélection des hyperparamètres sur la validation}
    (4\,886 exemples), avec le F1-score comme critère :
    \begin{itemize}
      \item seuil de décision (tous les modèles) : grille
        $\{0{,}3 ;\, 0{,}4 ;\, 0{,}5 ;\, 0{,}6\}$ ;
      \item profondeur maximale de l'arbre : grille $\{3, 5, 7, 9\}$ ;
      \item nombre de voisins du $k$-NN : grille $\{5, 11, 21, 31\}$ ;
    \end{itemize}
  \item \textbf{évaluation finale unique} sur les 4\,886 exemples du test ;
  \item \textbf{vérification} : chaque modèle est ré-entraîné avec son
    équivalent scikit-learn \cite{pedregosa2011sklearn} sur les mêmes données
    et avec des hyperparamètres équivalents, afin de valider la justesse de nos
    implémentations. Cette comparaison joue un rôle de contrôle qualité, pas de
    solution.
\end{enumerate}

Les grilles sont volontairement compactes : il s'agit de démontrer la démarche
de sélection sur validation, pas de mener une recherche exhaustive.
L'expérience complète est reproductible (graine fixe, une seule commande) et
s'exécute en une vingtaine de secondes.

\subsection{Métriques d'évaluation}

Toutes les métriques sont implémentées from scratch
(\texttt{src/metrics.py}). En notant TP, TN, FP, FN les vrais/faux
positifs/négatifs (la classe positive étant le défaut) :

\begin{align*}
  \text{accuracy} &= \frac{TP + TN}{TP + TN + FP + FN}
  & \text{précision} &= \frac{TP}{TP + FP} \\[4pt]
  \text{rappel} &= \frac{TP}{TP + FN}
  & F_1 &= \frac{2 \cdot \text{précision} \cdot \text{rappel}}
               {\text{précision} + \text{rappel}}
\end{align*}

Le choix de ces métriques découle directement du déséquilibre des classes :

\begin{itemize}
  \item l'\textbf{accuracy} donne une vue d'ensemble mais est insuffisante
    seule (prédire systématiquement la classe majoritaire donne déjà
    $78\,\%$) ;
  \item la \textbf{précision} mesure la fiabilité des alertes : parmi les
    défauts prédits, combien sont réels ;
  \item le \textbf{rappel} mesure la couverture : parmi les vrais défauts,
    combien sont détectés. C'est la métrique critique pour la banque,
    puisqu'un faux négatif correspond à un crédit accordé à un client qui ne
    remboursera pas ;
  \item le \textbf{F1-score}, moyenne harmonique des deux précédentes, sert de
    critère de sélection des hyperparamètres ;
  \item la \textbf{matrice de confusion} permet d'analyser la nature des
    erreurs ;
  \item la \textbf{courbe ROC} et l'\textbf{aire sous la courbe (AUC)}
    \cite{fawcett2006roc} comparent les modèles indépendamment du seuil de
    décision, à partir des probabilités prédites : l'AUC s'interprète comme la
    probabilité qu'un défaut tiré au hasard reçoive un score plus élevé qu'un
    non-défaut tiré au hasard.
\end{itemize}

Les métriques de régression (MAE, RMSE, $R^2$) et de partitionnement (inertie,
silhouette) ne sont pas pertinentes ici : le problème est une classification
supervisée, ni une régression ni un clustering.
